\begin{textsample}{POS Dim 2 – 2009\_c6 – Score 1.00 – 2009\_c6/enade\_2009\_com\_soc\_rp\_c9c58bdcc7.txt}  \label{ex:f2_pos_295}
Capítulo 3 Análise Técnica da Prova Este capítulo tem por objetivo apresentar o desempenho dos estudantes de COMUNICAÇÃO SOCIAL – RELAÇÕES PÚBLICAS no ENADE/2009.
Para isso, foram calculadas as estatísticas básicas da prova como um todo, bem como as estatísticas das questões relacionadas à Formação Geral e ao Componente Específico.
Nas Tabelas, são apresentadas as seguintes estatísticas : total da população e de presentes, média, erro-padrão da média, desvio-padrão, mediana, nota máxima, nota mínima e coeficiente de assimetria ( ver definição em 1.4.3 ).
Nas estatísticas das questões discursivas, foram detalhadas também as notas nulas : em branco, por protesto ou por insuficiência.
As estatísticas apresentadas neste capítulo contemplam, separadamente, os ingressantes, os concluintes e o total de estudantes, e foram calculadas tendo -se em vista as seguintes agregações : ( a ) Região Geográfica e País e ( b ) a Categoria Administrativa e a Organização Acadêmica.
As notas brutas das questões discursivas em cada uma das partes, Formação Geral/Componente Específico, foram calculadas a partir da média simples das notas de cada uma das questões correspondentes.
A primeira questão discursiva da Formação Geral foi desconsiderada e, portanto, a nota bruta referente ás questões discursivas de Formação Geral foi igual à nota da questão mantida, a segunda.
Na Componente Específico a nota bruta correspondeu à média das notas nas três questões.
As notas brutas das questões objetivas de cada uma das partes, Formação Geral/Componente Específico, foram calculadas a partir dos acertos dos estudantes considerando -se uma regra de três para a conversão para a escala de 0 a 100, descontando -se aquelas questões que foram desconsideradas pela comissão ou que não passaram pelo teste do índice de discriminação ponto-bisserial ( ver seção 1.4.6 para a descrição do índice ).
Questões com índices de discriminação inferior a 0,2 foram excluídas do cômputo da nota do estudante.
As notas brutas de Formação Geral e no Componente Específico para cada um dos estudantes foram calculadas como médias ponderadas das suas notas nas questões objetivas e discursivas, ponderação diferente para cada parte : C\_FG = C\_FG\textasciicircum{}OBJ × 0,6 + C\_FG\textasciicircum{}DIS × 0,4 C\_CE = C\_CE\textasciicircum{}OBJ × 0,85 + C\_CE\textasciicircum{}DIS × 0,15, onde C\_FG é a nota bruta da Formação Geral, C\_OBJ\_FG é a nota bruta nas questões objetivas da Formação Geral, C\_DIS\_FG é a nota bruta nas questões discursivas da Formação Geral, C\_CE é a nota bruta nas questões objetivas do Componente Específico, C\_OBJ\_CE é a nota bruta nas questões objetivas do Componente Específico, e C\_DIS\_CE é a nota bruta nas questões discursivas do Componente Específico.
A nota bruta da prova, C\_GER, por sua vez, é também uma média ponderada da nota de cada uma das partes, Formação Geral/Componente Específico : C\_GER = 0,25 × C\_FG + 0,75 × C\_CE. 3.1 Estatísticas Básicas da Prova 3.1.1 Estatísticas Básicas Gerais A Tabela 3.1 apresenta as estatísticas básicas da prova por grupo de estudantes.
O total de estudantes inscritos para a prova foi de 5.045.
Destes, 80,8 \% compareceram, sendo que a presença foi maior entre concluintes ( 87,5 \% ) do que entre ingressantes ( 73,3 \% ).
A média geral da prova foi 45,76, sendo que os ingressantes obtiveram média mais baixa, 43,19, que os concluintes, cuja média foi 47,69.
O desvio-padrão geral foi de 17,8, sendo o do grupo de ingressantes maior ( 17,8 ), do que o do grupo dos concluintes ( 17,5 ), aqueles possuíram uma distribuição menos homogênea nas notas.
A nota máxima foi 94,3, obtida por pelo menos um concluinte ao passo que a maior nota obtida por um ingressante foi 86,2.
Os coeficientes de assimetria para os dois grupos considerados, o de concluintes e o de ingressantes, bem como para a população como um todo foram todos negativos, indicando uma maior concentração à direita da distribuição e um maior espalhamento à esquerda.
Uma comparação entre o desempenho de ingressantes e concluintes pode ser realizada a partir dos dados apresentados no Gráfico 3.1.
Cumpre destacar que as notas dos concluintes, de um modo geral, foram mais elevadas do que as dos ingressantes, cuja distribuição de notas está à esquerda dos primeiros.
A ideia de melhor desempenho é reforçada pelas notas medianas dos dois grupos : 47,7 entre os concluintes e 43,2 entre os ingressantes.
Nas faixas de notas até 50 pontos o percentual de ingressantes é maior que o dos concluintes, situação esta que se inverte para as faixas de notas superiores a 50 pontos.
Ainda assim, aproximadamente 39,3 \% dos ingressantes situaram -se nas faixas de notas superiores a 50,0 pontos, por oposição a 52,1 \% dos concluintes.
O Gráfico 3.2, a seguir, apresenta informações referentes ao desempenho geral de ingressantes e concluintes, comparando os resultados considerando a Região, categoria administrativa e organização acadêmica.
Levando -se em conta as notas médias dos alunos em cada Região, observa -se que, nos dois grupos, há uma variação expressiva entre as Regiões.
A nota média mais elevada entre os concluintes ocorreu na Região Sul ( 53,25 ), acima da média nacional ( 47,68 ).
Em relação aos ingressantes, as médias mais altas foram obtidas nas Regiões Sul ( 46,22 ) e Centro-Oeste ( 45,18 ) nas demais Regiões a média foi inferior à média nacional ( 43,18 ).
Na Região Nordeste, foi obtida a média de notas mais baixa entre os ingressantes ( 41,58 ).
Entre os concluintes, a nota mais baixa foi verificada na Região Norte ( 40,57 ).
O Norte apresentou a média dos ingressantes superior a dos concluintes.
Este curso é oferecido em instituições federais, estaduais, municipais e particulares.
Levando -se em conta os agrupamentos dos estudantes em categorias administrativas, observa -se que a pontuação média mais elevada entre os concluintes e ingressantes foi encontrada em instituições de origem municipal ( respectivamente 55,63 e 46,80 ) – ambas acima da média nacional.
Tendo como foco as notas médias dos estudantes agrupados por organização acadêmica, notamos que universidades apresentam as maiores médias ( concluintes – 48,49 e ingressantes – 43,64 ).
Ao passo que as piores notas médias foram observadas nas Faculdades para concluintes ( 46,15 ) e Centros Universitários para ingressantes ( 40,76 ). 3.1.2 Estatísticas Básicas de Formação Geral A Tabela 3.2 apresenta as estatísticas básicas de Formação Geral por grupo de estudantes.
O total de estudantes inscritos para a prova foi de 5.045.
Destes, 80,8 \% compareceram, sendo que a presença foi maior entre concluintes ( 87,5 \% ) do que entre ingressantes ( 73,3 \% ).
A média geral da parte da prova de Formação Geral foi de 42,13, sendo que os ingressantes obtiveram média mais baixa, 40,90, que os concluintes, cuja média foi 43,04.
O desvio-padrão geral foi de 22,1 para ingressantes e concluintes.
A nota máxima foi 91,4, obtida por estudantes de ambos os grupos.
Para esta parte, os coeficientes de assimetria são positivos para o total e os ingressantes.
O Gráfico 3.3 permite a comparação entre concluintes e ingressantes relativamente ao desempenho na parte da prova que aborda a Formação Geral.
Nesta parte, se excluídas as provas com nota zero, tal como na prova como um todo, também os concluintes apresentaram desempenho superior, embora seja perceptível um maior equilíbrio.
A mediana dos concluintes foi igual à mediana geral, 42,8 enquanto que a dos ingressantes foi menor, 41,1.
Já a média dos concluintes é superior à dos ingressantes : 43,0 e 40,9, respectivamente.
Dos ingressantes, aproximadamente 35,0 \% alcançaram nota superior a 50 ; dentre os concluintes, o percentual dos estudantes nesta faixa de notas foi muito semelhante e igual a 38,2 \%.
No Gráfico 3.4, são apresentadas as informações referentes ao desempenho dos ingressantes e dos concluintes na parte de Formação Geral, em diferentes agregações como Região, categoria administrativa e organização acadêmica.
Considerando -se as notas médias dos estudantes segundo cada Região, observou -se que, tanto no grupo dos ingressantes quanto no grupo dos concluintes, houve diferenças marcantes.
As médias variam, entre os concluintes, de 30,39 na Região Norte a 50,58 na Região Sul.
Entre os ingressantes, os valores foram de 38,79 na Região Sudeste a 45,32 na Região Sul.
A maior dispersão de notas dos alunos foi encontrada na Centro-Oeste tanto para os concluintes como para os ingressantes.
Nas Regiões Sudeste e Sul, os concluintes apresentaram médias maiores do que os ingressantes, a situação \textbf{inversa} aconteceu no Norte, Nordeste e Centro-Oeste.
Todas as Regiões apresentaram concluintes.
Tendo como foco as categorias administrativas, destacou -se que todas as instituições estavam representadas, sendo que as notas médias das federais foram superiores para ingressantes ( 44,59 ).
E as notas médias das municipais foram superiores para concluintes ( 54,02 ).
Considerando -se o tipo de organização acadêmica, as Universidades apresentaram os maiores valores, tanto para concluintes ( 44,40 ) quanto para ingressantes ( 42,15 ).
Os alunos dos Centros Universitários se apresentam menos homogêneos, já que o desvio padrão é o maior encontrado em ambos os grupos.
Os Centros Federais de Educação Tecnológica e os Institutos Federais de Educação, Ciência e Tecnologia não estavam representadas nesta área. 3.1.3 Estatísticas Básicas do Componente Específico A Tabela 3.3 apresenta as estatísticas básicas referentes ao Componente Específico da prova, tendo sua média geral, de 47,02, superior ao desempenho da prova de COMUNICAÇÃO SOCIAL – RELAÇÕES PÚBLICAS como um todo, cuja média foi 45,76, como apresentado na Tabela 3.1, que por sua vez foi superior ao desempenho na que avalia na Formação Geral, que teve média igual a 42,13, conforme mostrado na Tabela 3.2.
Os concluintes apresentaram um desempenho médio ( 49,28 ) superior ao dos ingressantes ( 44,00 ), havendo ainda indícios de maior homogeneidade entre as notas médias dos concluintes do que entre as dos ingressantes, como sugeriu o desvio padrão, menor no grupo dos concluintes do que no dos ingressantes.
As notas máximas obtidas por concluintes ( 98,2 ) e ingressantes ( 89,4 ) foram diferentes.
As distribuições de notas para cada grupo se apresentaram com uma inclinação para notas maiores, conforme evidenciado pelos coeficientes de assimetria negativos.
A situação de assimetria é confirmada pelo Gráfico 3.5, que proporciona uma comparação do desempenho entre concluintes e ingressantes em relação ao Componente Específico da prova.
Claramente a distribuição das notas médias dos ingressantes, à esquerda daquela dos concluintes, evidenciou o percentual mais elevados de notas médias nas faixas mais baixas para este grupo.
As medianas das distribuições das notas dos ingressantes e dos concluintes foram 46,1 e 51,0, respectivamente, enfatizando o melhor desempenho dos concluintes nesta parte da prova.
Assim como os Gráficos 3.2 e 3.4, o Gráfico 3.6 apresenta uma comparação dos resultados em relação às Regiões, Categoria Administrativa e Organização Acadêmica, agora levando em conta o desempenho dos ingressantes e dos concluintes no Componente Específico da prova.
Tendo como foco as notas médias dos estudantes, em determinada Região, observa -se que, entre os concluintes, as mais elevadas foram encontradas na Região Sul ( 54,19 ), valores acima da média nacional para esta parte ( 49,28 ).
Entre os ingressantes, as maiores notas ocorreram nas Regiões Norte ( 44,08 ), Sul ( 46,56 ) e Centro-Oeste ( 46,90 ), todas acima da média nacional ( 43,99 ).
Novamente, para análise segundo a Categoria Administrativa, todas as entidades foram representadas.
Observa -se comportamento diverso entre a Formação Geral e o Componente Específico : na Formação Geral, as notas mais elevadas para concluintes estavam na categoria municipal e para ingressantes estavam na federal, já no Componente Específico, as notas mais elevadas, tanto para concluintes ( 56,21 ) quanto para ingressantes ( 49,56 ), estavam na origem municipal.
Os estudantes das instituições particulares apresentam uma homogeneidade, com desvio padrão das notas menores para os concluintes do que para os ingressantes.
Quanto à Organização Acadêmica, o comentário é o mesmo já feito para a nota bruta e referente à Formação Geral : somente Universidades, Centros Universitários e Faculdades são representados neste curso.
Considerando -se o tipo de Organização Acadêmica, entre os concluintes, Universidades apresentaram os maiores valores, seguidas dos Centros Universitários e das Faculdades.
E entre os ingressantes, Faculdades apresentaram os maiores valores, seguidas das Universidades e dos Centros Universitários. 3.2 Análise das Questões Objetivas 3.2.1 Formação Geral A Tabela 3.4 apresenta as estatísticas básicas relativas às oito questões objetivas da parte da prova que abrange a Formação Geral dos estudantes.
Adiantamos que a primeira questão foi desconsiderada pelo critério da correlação ponto-bisserial.
Como pode ser observado na Tabela 3.4, os estudantes concluintes obtiveram desempenho superior, com nota média de 50,5, em relação aos estudantes ingressantes cuja nota média foi 49,2.
A variabilidade entre os concluintes foi um pouco menor do que entre os ingressantes, como indicaram os desvios-padrão de 22,9, para o grupo de ingressantes e 21,7 para o grupo dos concluintes.
A amplitude total das notas foi a mesma para ambos os grupos : nota mínima de 0,00 e máxima de 85,7.
A mediana foi igual para ambos os grupos : 57,1.
A assimetria, medida pelo coeficiente de assimetria ( skewness ) foi também a mesma para ingressantes e os concluintes ( -0,5 ), uma assimetria para a direita, fato reforçado pelos valores menores das médias dos dois grupos quando comparados à mediana correspondente.
A Tabela 3.5 apresenta a classificação das questões objetivas da Formação Geral, segundo o índice de facilidade ( questões em vermelho foram desconsideradas pelo critério ponto-bisserial ).
Como as questões foram comuns a todas as áreas, a tabela apresenta a classificação considerando -se os concluintes de todas as áreas e aqueles do curso de COMUNICAÇÃO SOCIAL – RELAÇÕES PÚBLICAS.
De acordo com os índices obtidos, as questões objetivas da prova foram assim avaliadas para a área : das oito questões, três tiveram entre 60 \%, inclusive, e 85 \% de acertos e foram classificadas como fáceis ( somente uma para o total de concluintes ) ; outras duas foram consideradas medianas, alcançando entre 40 \%, inclusive, e 60 \% de acertos ( três para o total de concluintes ) ; duas obtiveram entre 15 \%, inclusive, e 40 \% de acertos, sendo consideradas difíceis ( três para o total de concluintes ).
Por fim, uma foi classificada como Muito difícil tento obtido percentual menor do que 15 \% de acertos ( uma para o total de concluintes ), justamente a questão eliminada pelo critério ponto-bisserial.
Nenhuma questão foi classificada na categoria Muito fácil.
A Tabela 3.6, descreve os resultados das análises das questões objetivas relativas à Formação Geral, segundo o poder de discriminação, utilizando -se para tal o índice de discriminação ponto-bisserial.
Como as questões foram comuns a todas as áreas, a tabela apresenta a classificação considerando -se os concluintes de todas as áreas e aqueles do curso de COMUNICAÇÃO SOCIAL – RELAÇÕES PÚBLICAS.
A desconsideração da questão segundo este critério considerou o universo de todos os concluintes.
Duas das oito questões apresentaram índices iguais ou acima de 0,40 ( nenhuma para a população de concluintes como um todo ), sendo, portanto, classificadas como muito boas nesta característica, para este grupo de estudantes, uma foi classificada como boa com índice entre 0,30, inclusive, e 0,40 ( três para os concluintes como um todo ) e quatro foram classificadas como médias, com índice entre 0,20, inclusive, e 0,30 ( quatro para os concluintes como um todo ).
Apenas uma das questões foi classificada como fraca para os concluintes de modo geral e também para os concluintes de COMUNICAÇÃO SOCIAL – RELAÇÕES PÚBLICAS, sendo desconsiderada pelo critério. uma para o total de concluintes ) ; outras duas foram consideradas medianas, alcançando entre 40 \%, inclusive, e 60 \% de acertos ( três para o total de concluintes ) ; duas obtiveram entre 15 \%, inclusive, e 40 \% de acertos, sendo consideradas difíceis ( três para o total de concluintes ).
Por fim, uma foi classificada como Muito difícil tento obtido percentual menor do que 15 \% de acertos ( uma para o total de concluintes ), justamente a questão eliminada pelo critério ponto-bisserial.
Nenhuma questão foi classificada na categoria Muito fácil.
A Tabela 3.6, descreve os resultados das análises das questões objetivas relativas à Formação Geral, segundo o poder de discriminação, utilizando -se para tal o índice de discriminação ponto-bisserial.
Como as questões foram comuns a todas as áreas, a tabela apresenta a classificação considerando -se os concluintes de todas as áreas e aqueles do curso de COMUNICAÇÃO SOCIAL – RELAÇÕES PÚBLICAS.
A desconsideração da questão segundo este critério considerou o universo de todos os concluintes.
Duas das oito questões apresentaram índices iguais ou acima de 0,40 ( nenhuma para a população de concluintes como um todo ), sendo, portanto, classificadas como muito boas nesta característica, para este grupo de estudantes, uma foi classificada como boa com índice entre 0,30, inclusive, e 0,40 ( três para os concluintes como um todo ) e quatro foram classificadas como médias, com índice entre 0,20, inclusive, e 0,30 ( quatro para os concluintes como um todo ).
Apenas uma das questões foi classificada como fraca para os concluintes de modo geral e também para os concluintes de COMUNICAÇÃO SOCIAL – RELAÇÕES PÚBLICAS, sendo desconsiderada pelo critério.
Na sequência, no Quadro 3.1, estão especificadas, por questão, as habilidades aferidas pelas questões de Formação Geral da prova da COMUNICAÇÃO SOCIAL – RELAÇÕES PÚBLICAS, bem como os resultados da análise relativa aos índices de facilidade e de discriminação de cada questão.
O índice de facilidade teve uma variação considerável, oscilando entre 0,10 até 0,80 ( na população de concluintes o intervalo foi de 0,08 a 0,75 ) e o de discriminação ficou entre 0,03 e 0,49 ( na população de concluintes o intervalo foi de 0,02 a 0,39 ).
A questão 8 aferiu as habilidades de “ Ler e interpretar textos ”, “ Analisar e criticar informações ”, “ Extrair conclusões por indução e/ou dedução ”, “ Estabelecer relações, comparações e contrastes em diferentes situações ”, “ Detectar contradições ”, “ Fazer escolhas valorativas avaliando consequências ”, “ Questionar a realidade ”, e “ Argumentar coerentemente ” foi a que apresentou, para a população de concluintes de COMUNICAÇÃO SOCIAL – RELAÇÕES PÚBLICAS, maior poder discriminatório com índice de 0,49.
A questão mais fácil, com 80 \% de acertos, foi também a de número 8.
De outra parte, a questão com menor índice de facilidade e de discriminação deste conjunto foi a questão 1, que foi desconsiderada, e que deveria aferir as mesmas habilidades da questão 8.
O Gráfico 3.7 mostra, a título de ilustração, o comportamento da questão 7 de Formação Geral para os estudantes concluintes da área de COMUNICAÇÃO SOCIAL – RELAÇÕES PÚBLICAS como função da geral da prova.
O poder de discriminação foi muito bom para os estudantes da área e bom para o conjunto de todos os concluintes.
Neste Gráfico, cada uma das cinco curvas, representou o percentual de respostas em determinada alternativa da questão, em função do escore total destes estudantes.
A curva azul marinho corresponde à alternativa A, a correta para esta questão.
Assim, observa -se que entre os estudantes com escores mais baixos – desempenho geral mais fraco na prova como um todo – além da alternativa correta, houve a escolha de alternativas incorretas, especialmente a alternativa B – curva verde claro seguida da alternativa E – curva em roxo.
A proporção de estudantes que selecionaram a resposta correta A, aumentou com o desempenho deste grupo de estudantes chegando ao acerto de 100 \% dos estudantes com notas pelo menos 87,5.
Esta análise permitiu verificar como a questão discriminou entre os dois grupos de desempenho de acordo com o índice obtido.
Os gráficos relativos às demais questões de Formação Geral constam do Anexo I. 3.2.2 Componente Específico A Tabela 3.7 apresenta as estatísticas básicas em relação às questões objetivas do Componente Específico.
Como pode ser observado, os estudantes concluintes obtiveram desempenho superior, mas por uma pequena margem com uma média de 52,13 pontos, em relação aos estudantes ingressantes com uma média de 46,65 pontos.
Os desvios-padrão indicaram que a variabilidade entre os concluintes foi menor ( 19,0 ) do que entre os ingressantes ( 19,5 ).
Os coeficientes de assimetria para os dois grupos são negativos indicando uma inclinação para a direita.
A assimetria foi, porém, maior em módulo entre os concluintes.
A análise das questões objetivas do Componente Específico em COMUNICAÇÃO SOCIAL – RELAÇÕES PÚBLICAS, em relação ao grau de facilidade da prova ( ver Tabela 3.8 ), mostrou que, das 27 questões, uma encontram -se na categoria Muito difícil ( abaixo de 15 \% de acertos ), seis são de nível Difícil ( entre 15 \% e 40 \% de acertos ), nove estão na categoria médio ( entre 40 \% e 60 \% de acertos ) e sete tiveram a classificação Fácil ( entre 60 \% e 85 \% de acertos ).
Nenhuma questão foi classificada como Muito fácil.
Quatro questões foram desconsideradas pela Comissão Assessora da Área de COMUNICAÇÃO SOCIAL – RELAÇÕES PÚBLICAS.
Das 27 questões objetivas relativas ao Componente Específico ( ver Tabela 3.9 ), doze obtiveram o índice de discriminação ' Muito bom ', cinco obtiveram o índice de discriminação ' Bom ' e três tiveram nível de discriminação Médio.
O índice ' Fraco ' que eliminaria a questão foi alcançado por três questões.
Na sequência, no Quadro 3.2 estão especificadas, por questão, as habilidades aferidas pelas questões do Componente Específico da prova, bem como os resultados da análise relativa aos índices de facilidade e de discriminação de cada questão.
Das 27 questões objetivas relativas ao Componente Específico ( ver Tabela 3.9 ), doze obtiveram o índice de discriminação ' Muito bom ', cinco obtiveram o índice de discriminação ' Bom ' e três tiveram nível de discriminação Médio.
O índice ' Fraco ' que eliminaria a questão foi alcançado por três questões.
Na sequência, no Quadro 3.2 estão especificadas, por questão, as habilidades aferidas pelas questões do Componente Específico da prova, bem como os resultados da análise relativa aos índices de facilidade e de discriminação de cada questão.
O Gráfico 3.8 mostra, a título de ilustração, o comportamento da questão 11 do Componente Específico para os estudantes concluintes da área de COMUNICAÇÃO SOCIAL – RELAÇÕES PÚBLICAS com função da nota geral na prova5.
Trata -se de uma questão considerada de dificuldade média para os concluintes de COMUNICAÇÃO SOCIAL – RELAÇÕES PÚBLICAS.
O poder de discriminação foi considerado muito bom para os estudantes da área.
Neste Gráfico, cada uma das cinco curvas, representa o percentual de respostas em determinada alternativa da questão, em função da nota total dos estudantes concluintes de COMUNICAÇÃO SOCIAL – RELAÇÕES PÚBLICAS.
A curva verde claro correspondeu à alternativa B, a correta para esta questão.
Assim, observou -se que entre os estudantes com escores mais baixos – desempenho geral mais fraco na prova como um todo –, houve a escolha de alternativas incorretas como a alternativa E – curva em roxo.
A proporção de estudantes que selecionaram a resposta correta aumenta com o desempenho deste grupo de estudantes chegando ao acerto de um pouco mais de 90 \% dos estudantes com notas no grupo [ 75,0 ; 87,5 ) caindo no grupo seguinte.
Esta análise permitiu verificar como a questão discriminou entre os dois grupos de desempenho de acordo com o índice obtido.
Os Gráficos relativos às demais questões do Componente Específico constam do Anexo I. 3.3 Análise das Questões Discursivas 3.3.1 Formação Geral A Tabela 3.10 apresenta as estatísticas básicas da prova por grupo de estudantes.
A análise dos resultados de desempenho dos estudantes nas questões discursivas que são apresentados na tabela 3.10 e no gráfico 3.9 mostrou que as notas foram mais baixas no conjunto dessas questões que no das questões objetivas.
O total de estudantes inscritos para a prova foi de 5.045.
Destes, 80,8 \% compareceram, sendo que a presença foi maior entre concluintes ( 87,5 \% ) do que entre ingressantes ( 73,3 \% ).
A média geral da prova foi de 30,56, sendo que os ingressantes obtiveram média mais baixa, de 28,59, que os concluintes, cuja média foi 32,03.
O desvio-padrão geral foi de 36,3, sendo o do grupo de ingressantes menor ( 34,8 ), indicando que, em contraposição ao grupo dos concluintes ( 37,3 ), aqueles possuíram uma distribuição mais homogênea nas notas.
A nota máxima foi 100,0, atingida por concluintes e ingressantes.
Os coeficientes de assimetria para os dois grupos considerados, o de concluintes e o de ingressantes, bem como para a população como um todo são todos positivos, indicando uma maior concentração à esquerda da distribuição.
Para todos os grupos, um pouco mais de metade dos estudantes presentes obtiveram nota zero nas questões discursivas da Formação Geral.
Das 2.238 provas com nota zero, 89 ( 4,0 \% ) continham um protesto, 852 ( 38,1 \% ) estavam em branco e 1.297 ( 58,0 \% ) foram respondidas, porém de forma inadequada.
Entre os concluintes, a maioria ( 58,8 \% ) dentre aqueles com nota zero apresentou com as provas respondidas de forma inadequada.
Entre os ingressantes, 56,8 \% das notas nulas decorreram de correção.
Uma comparação entre o desempenho de ingressantes e concluintes pode ser realizada a partir dos dados apresentados no Gráfico 3.9.
Cumpre destacar que as notas dos concluintes, de um modo geral, foram ligeiramente mais elevadas do que as dos ingressantes, cuja distribuição de notas está à esquerda dos primeiros.
A ideia de melhor desempenho foi reforçada pela proporção de notas nulas : 52,7 \% entre os concluintes e 54,1 \% entre os ingressantes.
Na sequência, discorrer -se-á, especificamente, sobre uma das questões discursivas de Formação Geral, a número 10 já que a número 9 foi desconsiderada.
Em especial, serão analisados os conteúdos englobados na questão e os resultados obtidos pelos estudantes.
Ao final, são apresentados os comentários da banca corretora a respeito do observado nas respostas dos estudantes, suas impressões e conclusões.
Cumpre esclarecer que, tendo em vista que as questões discursivas de Formação Geral foram padronizadas, ou seja, constam de todas as provas, os comentários da banca foram os mesmos para todas as áreas acadêmicas, eis que direcionados a todos os estudantes que participaram do ENADE/2009. 3.3.1.1 Análise da Questão Discursiva 9 de Formação Geral DESCONSIDERADA 3.3.1.2 Análise da Questão Discursiva 10 de Formação Geral A Questão Discursiva nº. 10 buscou avaliar as habilidades : • Ler e interpretar textos ; • Analisar e criticar informações ; • Extrair conclusões por indução e/ou dedução ; • Estabelecer relações, comparações e contrastes em diferentes situações ; • Fazer escolhas valorativas avaliando consequências ; • Questionar a realidade ; • Argumentar coerentemente.
O Quadro 3.3 apresenta as habilidades/competências, temas predominantes e estatísticas básicas da questão 10 de Formação Geral por grupo de estudantes.
Cumpre notar que como a questão 9 foi desconsiderada, a nota referente às questões discursivas da Formação Geral é idêntica à nota da questão 10.
A média geral desta questão foi de 30,56, sendo que os ingressantes obtiveram média mais baixa, de 28,59, que os concluintes, cuja média foi 32,03.
O desvio-padrão geral foi de 36,3, sendo o do grupo de ingressantes menor ( 34,8 ), indicando que, em contraposição ao grupo dos concluintes ( 37,3 ), aqueles possuíram uma distribuição mais homogênea nas notas.
A nota máxima foi 100,0, atingida por concluintes e ingressantes.
Os coeficientes de assimetria para os dois grupos considerados, o de concluintes e o de ingressantes, bem como para a população como um todo foram todos positivos, indicando uma maior concentração à esquerda da distribuição.
Para todos os grupos, um pouco mais de metade dos estudantes presentes tiveram nota zero nas questões discursivas da Formação Geral.
Das 2.238 provas com nota zero, 89 ( 4,0 \% ) continham um protesto, 852 ( 38,1 \% ) estavam em branco e 1.297 ( 58,0 \% ) foram respondidas, porém de forma inadequada.
Entre os ingressantes, a maioria ( 58,8 \% ) dentre aqueles com nota zero estavam com as provas respondidas de forma inadequada.
Entre os ingressantes, 56,8 \% das notas nulas decorreram do desempenho do próprio estudante verificadas na correção.
Uma comparação entre o desempenho de ingressantes e concluintes para esta questão foi exatamente igual à já realizada para a nota discursiva de Formação Geral a partir dos dados apresentados no Gráfico 3.9. 3.3.1.3 Comentários sobre as respostas à Questão Discursiva 10 Apresentação A questão 10 de Formação Geral apresentou um questionamento capaz de analisar a habilidade do estudante em reconhecer o que é um direito social, optar por um direito social vigente que possa diminuir as desigualdades de renda familiar no Brasil e discutir os argumentos em favor do direito social escolhido como suporte à resposta.
Para análise das respostas apresentadas pelos estudantes, houve uma leitura inicial bastante criteriosa.
Foram considerados para efeito de pontuação, os seguintes argumentos : • Acesso à educação pública, gratuita e de qualidade, garantindo ao cidadão melhores condições de vida ; • Assistência à saúde, em seu contexto mais amplo, garantindo os cuidados essenciais de saúde, baseados em técnicas apropriadas ao alcance de todas as pessoas, contando com a participação da população ; • Permanência do estudante na escola, em todos os níveis escolares ( da educação infantil à educação superior, possibilitando a qualificação profissional ), o que garante melhores condições de trabalho e remuneração.
Considerou -se a qualificação e o desempenho docente, desenvolvendo nos estudantes melhores expectativas de aprendizagem e interesse pelo conteúdo da escola.
Os professores desenvolvendo metodologias adequadas e conhecimento pleno quanto a questões de aprendizagem ; • Condições dignas de trabalho, com remuneração que garanta qualidade de vida.
Foram consideradas as condições de trabalho, a preparação profissional para o trabalho e os cursos de aperfeiçoamento, oportunizando momentos de aprendizagem e desenvolvimento, a partir de situações desencadeadoras ; • Ser proprietário do imóvel em que reside, reduzindo os gastos com aluguel, garantindo melhores condições de vida, equilíbrio financeiro e construção do bem-estar familiar.
Desempenho dos estudantes na questão É importante salientar que havia respostas que desenvolveram adequadamente o questionamento apresentado, com fundamentação baseada em argumentos coerentes.
Contudo, em um número considerável de respostas, havia apenas um argumento e não dois, conforme fora solicitado.
Muitos estudantes informaram sobre problemas da educação, de um modo geral.
Algumas observações mencionadas pelos corretores durante o processo de correção foram : • há preocupação com uma educação de qualidade, com sugestões de cursos profissionalizantes e métodos preparatórios de aulas. • a educação é revelada como de baixa qualidade, exaltando -se treinamentos e cursos de atualização para a classe docente. • a saúde como essência para a vida.
O SUS é colocado como parte integrante de direito social, garantindo às pessoas fácil acesso. • oportunidade de emprego como forma de ascensão social, direito, cidadania. • os políticos brasileiros como causadores do despreparo educacional.
As respostas analisadas apresentaram norma culta, coesão e coerência razoáveis, que garantiram a compreensão da mensagem desenvolvida pelos estudantes.
Os erros ortográficos foram poucos e não comprometeram a discussão sobre o assunto e sobre os pontos de vista abordados.
Quanto à coerência, houve um desvio quanto ao objetivo da questão, já que alguns estudantes, através dos seus argumentos retrataram reivindicações, reclamações e situações de problemas enfrentados pelos brasileiros, citando golpe políticos.
Houve erros quanto á concordância verbal e nominal.
Houve, também, dificuldade em utilizar os conectores corretamente.
Alguns estudantes, apresentaram em suas respostas redundância e prolixidade.
O desempenho foi considerado satisfatório, pois a maioria dos estudantes soube responder ao questionamento.
Houve domínio do conteúdo, pois as ideias, os argumentos e a originalidade corresponderam ao esperado. 3.3.2 Componente Específico Nas questões discursivas do Componente Específico – questões de números 38, 39 e 40 –, a diferença de desempenho entre os grupos não foi muito acentuada, inclusive menor do que para as questões discursivas de Formação Geral.
A Tabela 3.11 contém as estatísticas da média destas questões.
Observa -se que no grupo dos concluintes a nota média alcançada foi 33,47, enquanto que no grupo dos ingressantes foi de 29,27, ambas maiores do que as obtidas nas questões discursivas de Formação Geral, que foram de 28,59 para os ingressantes e 32,03 para os concluintes ( Tabela 3.10 ).
As médias foram também menores dos que as obtidas para as questões objetivas do Componente Específico ( Tabela 3.7 ) : respectivamente 46,65 e 52,13 para ingressantes e concluintes.
A análise de cada uma destas questões será feita nas subseções seguintes.
Os valores positivos para a medida de assimetria indicaram uma distribuição com cauda para a esquerda.
Dos 4.076 presentes, 27,3 \% obtiveram nota zero.
A tabela 3.12 apresenta um detalhamento da situação destes estudantes.
A proporção dos que não fizeram nenhuma das 3 questões foi de 61,5 \%.
Esta proporção foi maior entre os ingressantes ( 65,2 \% ) do que entre concluintes ( 58,5 \% ).
A proporção de provas com todas as questões respondidas correspondeu a 6,2 \% das provas.
Esta proporção foi maior entre os concluintes ( 7,1 \% ) do que entre os ingressantes ( 4,9 \% ).
Houve um total de 15 provas ( 1,4 \% ) com as 3 questões com protesto.
O Gráfico 3.10 apresenta a distribuição das notas nas questões discursivas do Componente Específico por grupo de estudantes.
A moda se situou no primeiro intervalo, aquele com valores entre 0 e 10 incluindo os extremos do intervalo.
A impressão de assimetria informada pelos coeficientes de assimetria positivos ( Tabela 3.11 ) foi reforçada no gráfico pela concentração no primeiro intervalo. 3.3.2.1 Análise da Questão Discursiva 38 do Componente Específico A questão 38, cujas habilidades aferidas encontram -se descritas no Quadro 3.4, que ainda contém outras informações, seguido da Tabela 3.13 e do Gráfico 3.11, foi aquela em que os grupos de estudantes observados obtiveram o melhor desempenho entre as questões discursivas do Componente Específico.
A nota média do grupo dos ingressantes foi de 38,63, e a dos concluintes 43,67.
O percentual de notas zero obtidas nos dois grupos, o mais baixo entre as questões discursivas do Componente Específico, foi de 33,7 \%, maior entre os ingressantes ( 35,2 \% ) do que entre os concluintes ( 32,5 \% ).
Quase 80 \% dos concluintes e dos ingressantes que tiraram zero deixaram a questão em branco ( ver Tabela 3.13 ).
Uma pequena minoria ( 7,6 \% ) utilizou a questão para protestar.
Apesar do baixo desempenho, em ambos os grupos a nota máxima de 100,0 pontos foi alcançada.
Os coeficientes de assimetria positivos indicam uma inclinação para a esquerda e cauda para a direita.
O Gráfico 3.11 apresenta a distribuição das notas na questão 38 discursiva do Componente Específico por grupo de estudantes.
A moda se situou no primeiro intervalo, aquele com valores entre 0 e 10 incluindo os extremos do intervalo.
A impressão de assimetria informada pelos coeficientes de assimetria positivos ( Quadro 3.4 ) foi reforçada no gráfico pela concentração no primeiro intervalo.
A distribuição apresentou outros máximos locais no sétimo, no nono e no último intervalo. 3.3.2.2 Comentários sobre as respostas à Questão Discursiva 38 A questão, de dificuldade mediana, tratou -se de tema atual e relevante às organizações modernas que é a responsabilidade social, suas ações e resultados.
O objetivo central da questão era apurar o conhecimento dos estudantes sobre as possibilidades no relacionamento empresa/responsabilidade social.
Com um enunciado longo, mas extremamente bem construído, a questão levou o estudante a raciocinar de forma objetiva e tópica, pois demandou apenas a criação de duas estratégias focadas na responsabilidade social e suas justificativas.
A correção foi facilitada pelo padrão de resposta que definiu de forma clara e objetiva, os elementos fundamentais à proposição das ações, não restringindo as possibilidades de temas propostos.
Desta forma, foi possível fazer uma ampla avaliação das respostas apresentadas e considerar todas que se enquadraram nos fundamentos exigidos.
Boa parte dos estudantes conseguiu produzir em suas respostas um conteúdo que satisfez em parte o comando da questão.
Não foi significativo o contingente de respostas com grau máximo, mas muitos estudantes tiveram um desempenho bom ou razoável, atingindo com as seguintes situações : sendo frequentes as respostas em que apenas foram estabelecidas as estratégias, esquecendo os estudantes das justificativas ; outras respostas em que apenas uma estratégia era criada e justificada, faltando a segunda ação solicitada pelo enunciado ; outras respostas em que era concebida uma só estratégia reelaborada na segunda metade da resposta, ficando patente a tentativa do estudante de obter pontuação dupla pela mesma solução ; outras respostas em que uma única justificativa, em geral mais ligada aos aspectos filosóficos da responsabilidade social era aplicada a duas estratégias para as quais só era aplicável como norma ética de cunho mais geral, e não como justificativa efetiva da ação.
Concluindo, podemos afirmar que o enfoque foi de agrado dos estudantes e não foram registrados impropérios ou desabafos, quando da resposta a questão. 3.3.2.3 Análise da Questão Discursiva 39 do Componente Específico A questão 39, cujas habilidades aferidas encontram -se descritas no Quadro 3.5, que ainda contém outras informações, seguido da Tabela 3.14 e do Gráfico 3.12, foi aquela em que os grupos de estudantes observados obtiveram o desempenho intermediário entre as questões discursivas do Componente Específico.
A nota média do grupo dos ingressantes foi de 25,79, e a dos concluintes 27,80.
O percentual de notas zero obtidas nos dois grupos foi de 49,3 \% entre os ingressantes e 46,1 \% entre os concluintes.
Das provas com notas nulas nesta questão, quase a metade foi deixada em branco : 49,3 \% entre os ingressantes e 43,1 \% entre os concluintes ( ver Tabela 3.14 ).
Entre os com nota nula, uma pequena minoria ( 2,4 \% ) utilizou a questão para protestar.
Apesar do baixo desempenho, em ambos os grupos a nota máxima de 100,0 pontos foi alcançada.
Os coeficientes de assimetria positivos indicaram uma inclinação para a esquerda e cauda para a direita.
O Gráfico 3.12 apresenta a distribuição das notas na questão 39 discursiva do Componente Específico por grupo de estudantes.
A moda se situou no primeiro intervalo, aquele com valores entre 0 e 10 incluindo os extremos do intervalo.
A impressão de assimetria informada pelos coeficientes de assimetria positivos ( Quadro 3.5 ) foi reforçada no gráfico pela concentração no primeiro intervalo. 3.3.2.4 Comentários sobre as respostas à Questão Discursiva 39 Questão de dificuldade mediana, que exigiu dos estudantes conhecimentos técnicos sobre públicos e técnicas de eventos, além de grande capacidade de analisar cuidadosamente as demandas propostas.
O objetivo central da questão era criar uma relação de quatro atividades de atendimento e recepção de convidados em um evento para a terceira idade.
O enunciado era claro e não propiciou qualquer tipo de confusão, ficando o estudante limitado apenas pelo seu conhecimento e bom senso.
O principal fator para a correção da resposta e um bom desempenho era a percepção das necessidades dos idosos e o respeito necessário no atendimento deste tipo de público.
A chave de resposta demandou a clara preocupação com o pós-evento, que deveria ser manifestada através de uma das ações propostas.
O equívoco mais frequente era provocado pelo esquecimento da solicitação de preocupação com o pós-evento.
No entanto, a listagem de ações de forma extremamente objetiva, ao ponto do entendimento da resposta ficar comprometido também foi frequente.
Por mais que o enunciado direcionasse a resposta para soluções tópicas, que ocupassem quatro ou cinco linhas, em muitos casos, as frases muito curtas e absolutamente sintéticas acabaram por reduzir a pontuação obtida pelos estudantes.
Detectou -se também alguma dificuldade, por parte dos estudantes, em \textbf{ordenar} adequadamente as ações propostas e, principalmente, em restringi -las ao “ atendimento ao público ”, como foi solicitado no enunciado, preferindo tratar do evento como um todo.
De uma forma geral, esta foi a questão de melhor desempenho dos estudantes entre as provas respondidas da habilitação. 3.3.2.5 Análise da Questão Discursiva 40 do Componente Específico A questão 40, cujas habilidades aferidas encontram -se descritas no Quadro 3.6, que ainda contém outras informações, seguido da Tabela 3.15 e do Gráfico 3.13, foi aquela em que os grupos de estudantes observados obtiveram o pior desempenho entre as questões discursivas do Componente Específico.
A nota média do grupo dos ingressantes foi de 23,48, e a dos concluintes 29,01.
Nesta questão, ocorre também o maior percentual de notas nulas : 54,2 \% entre os ingressantes e 49,9 \% entre os concluintes.
Aqui também a maioria das provas com nota zero decorreu de provas deixadas em branco e com a questão mal respondida, quase 50 \% entre ingressantes e concluintes ( ver Tabela 3.15 ).
Uma pequena minoria ( 2,7 \% ), dentre aqueles com nota nula, utilizou a questão para protestar.
Em ambos os grupos a nota máxima de 100,0 pontos foi alcançada.
Os coeficientes de assimetria positivos indicaram uma inclinação para a esquerda.
O Gráfico 3.13 apresenta a distribuição das notas na questão 40 discursiva do Componente Específico por grupo de estudantes.
A moda se situou no primeiro intervalo, aquele com valores entre 0 e 10 incluindo os extremos do intervalo.
A impressão de assimetria informada pelos coeficientes de assimetria positivos ( Quadro 3.6 ) foi reforçada no gráfico pela concentração no primeiro intervalo. 3.3.2.6 Comentários sobre as respostas à Questão Discursiva 40 Esta é uma das questões mais complexas apresentadas na área de Comunicação Social, abordando um tema atualíssimo e de grande importância para a formação do profissional de relações públicas.
Seu padrão de resposta exigiu dos estudantes amplo conhecimento pragmático, principalmente ao que tange à Comunicação Pública e a seus conceitos voltados à cidadania.
O objetivo central da questão era apurar o conhecimento dos estudantes sobre o novo paradigma da Comunicação, com as Relações Públicas.
O enunciado exigiu do estudante uma boa capacidade de compreensão do processo de comunicação como um todo, direcionou a resposta para ações de responsabilidade social e ambiental e de foco no cidadão e não no consumidor.
Esta foi a questão com menor aproveitamento entre as respostas dessa habilitação.
Houve grande dificuldade dos estudantes em perceber o que era solicitado pelo enunciado, evidenciando o baixo conhecimento existente em muitas instituições de ensino sobre a questão abordada.
O erro mais frequente foi a apresentação de respostas, cujo conteúdo era a mera \textbf{reprodução} do enunciado sem qualquer tipo de reflexão, ou o apelo ao comportamento ético dos profissionais.
Outro equívoco constante foi a explanação sobre a importância do profissional de Relações Públicas para a boa comunicação da empresa com seu público.
Também foi frequente a tentativa de “ impressionar ” a banca de correção, recheando o texto de anglicismos, termos técnicos e explicações entusiasmadas sobre a necessidade das empresas de deixarem a cargo de profissionais habilitados as atividades específicas de Relações Públicas.

% matched lemmas: inverso, ordenar, reprodução
\end{textsample}
